\chapter{To Smooth or Not to Smooth}
\label{ch:smoothing}

``The data is noisy --- let me smooth it'' is among the most natural
instincts in time-series work, and in this genre it is usually wrong in an
instructive way. This chapter builds a small theory of when temporal
filtering helps, when it destroys the only signal you have, and what the
correct operation is when the data arrives \emph{already} smoothed.

\section{Filters, in one page}

A causal linear filter maps a series $x_t$ to
$\tilde x_t = \sum_{j \ge 0} h_j\, x_{t-j}$. The two workhorses:

\begin{itemize}
\item \textbf{SMA($w$)}: $h_j = 1/w$ for $j < w$. Frequency response is the
  Dirichlet kernel of \S\ref{sec:fourier}: passes slow oscillations,
  crushes anything faster than $\sim 1/w$, with hard nulls at multiples of
  $1/w$. Group delay $\approx (w{-}1)/2$: \emph{the output lags the input
  by half a window}.
\item \textbf{EWMA($\alpha$)}: $h_j = \alpha(1-\alpha)^j$. Same low-pass
  character with an exponential window; effective span
  $\approx 2/\alpha - 1$; delay $\approx (1-\alpha)/\alpha$ steps. Prefer
  it in streaming systems ($O(1)$ state) and when you want no hard window
  edge.
\end{itemize}

Both are \textbf{low-pass filters}: they keep the slow component and delete
the fast one, \emph{and they delay what they keep}. Every use of smoothing
must answer two questions: \emph{is my signal slow or fast?} and \emph{can I
afford the delay?}

\section{Why smoothing the target is a category error here}

Chapter~\ref{ch:targets} showed the target already \emph{is} an SMA ---
the host applied the smoothing, and that construction is precisely what
makes the problem hard: averaging 20 future innovations pools their
unpredictability into one number whose conditional mean is nearly zero.
Additional smoothing of predictions or labels cannot help --- the metric
compares you to the raw stored target, and any filter you apply to
$\hat y$ moves it away from the conditional mean the metric rewards
(Chapter~\ref{ch:metric}). If you feel the urge to smooth your predictions,
what you have actually detected is \emph{prediction variance}, and the
correct treatments are shrinkage (calibrated, \S\ref{ch:metric}) and
ensembling (Chapter~\ref{ch:ensembling}) --- not a moving average over time.

\section{Smoothing inputs: the level/innovation decomposition}

For \emph{features}, the right mental model is decomposition. Any series
splits as
\[
x_t \;=\; \underbrace{\tilde x_t}_{\text{level (slow)}}
\;+\; \underbrace{(x_t - \tilde x_t)}_{\text{innovation (fast)}},
\]
where $\tilde x$ is a trailing smooth. The two components carry different
kinds of information and age at different speeds:

\begin{itemize}
\item The \textbf{level} is context: what regime is this symbol in, how
  volatile, how high relative to its own history. It changes slowly, so its
  half-window delay is harmless.
\item The \textbf{innovation} is timing: what just changed. It is fast, and
  a smoother \emph{delays and attenuates exactly it}.
\end{itemize}

\begin{keybox}[: give the model both parts; never let a smoother eat a
timing signal]
The correct engineering move is almost never ``smooth the feature'' but
``split the feature'': emit the trailing mean (or the raw level) \emph{and}
the deviation from it. The reference solution's strongest engineered block
does exactly this --- \code{feature $-$ rolling\_mean$_{1000}$} per symbol,
i.e.\ it keeps the \emph{high-pass} part and lets the slow level be handled
by normalization. Notice the direction: the naive instinct smooths;
the winning pipeline \emph{un}-smooths.
\end{keybox}

The atlas of Chapter~\ref{ch:featforensics} makes the warning quantitative.
The lead features \feat{37}/\feat{38} peak at $k = +3$: their information
concerns the next few steps. An SMA-20 applied to them would (i) delay
that information by $\sim$10 steps --- past its expiry --- and (ii) dilute
it 20-fold with stale values. One glance at a feature's lead--lag curve
tells you whether it can survive smoothing: \textbf{if the curve's mass is
at small $|k|$, smoothing kills it; if the mass is broad and slow,
smoothing merely denoises it.}

Where smoothing inputs \emph{does} pay: constructing deliberate multi-scale
representations. A fast feature accompanied by its EWMA-30 (both provided!)
lets a tree or network compute ``current value relative to its recent
norm'' --- which is the level/innovation split again, arrived at from the
other side.

\section{When the data arrives pre-smoothed: deconvolve}

The genre's sharpest smoothing insight is the inverse one. Forensics found
three features (\feat{12}/\feat{67}/\feat{70}) that are themselves SMA-160
of some hidden series (\S\ref{sec:featdeconv}) --- the host shipped them
already low-passed. For such features the informative part --- the recent
increment of the underlying series --- is attenuated by a factor
$\sim w$ and buried under a huge stale level. Three escalating remedies:

\begin{enumerate}
\item \textbf{Short difference}: $x_t - x_{t-\delta}$ removes the common
  stale mass; for an SMA-$w$ input, the difference over $\delta$ steps
  equals the mean of the $\delta$ newest underlying values minus the
  $\delta$ oldest --- a band-pass around ``what changed recently.'' Cheap,
  robust, first thing to try.
\item \textbf{Difference at the detected window}: $x_t - x_{t-w}$
  telescopes an SMA-$w$ into the mean of the $w$ newest underlying values
  minus the mean of the $w$ before them --- a clean two-window momentum
  with no stale overlap. Requires knowing $w$, which the ACF ramp gives
  you.
\item \textbf{Ridge deconvolution} (\S\ref{sec:deconv}): reconstruct the
  underlying series itself. Maximal resolution; requires the regularization
  machinery and validation checks; reserve it for series important enough
  to justify the care (for us: the responder paths, where the deconvolved
  $\hat s$ became the reference signal for the entire atlas).
\end{enumerate}

The derivation you were promised in Chapter~\ref{ch:targets}'s exercise
lands here: the first difference of SMA-$w$-of-white-noise has
autocorrelation $-1/2$ at... no --- work it out; the point of the exercise
is that differencing a smoothed series produces a specific, recognizable
negative-correlation signature at the window length, which is one more way
to \emph{confirm} a detected window before you build features on it.

\section{Smoothing across the cross-section}

One more direction exists in a panel: smoothing across \emph{symbols} at a
fixed time. The per-timestamp market average is exactly that --- a
cross-sectional smoother --- and unlike temporal smoothing it involves
\textbf{no delay}: all inputs are contemporaneous and legal. This asymmetry
is worth internalizing: in a streaming panel, \emph{the cross-section is
the one direction you can average without paying latency}. It is no
coincidence that market-average features appear in every strong solution
to this competition, while temporal smoothing of features appears in none
of them.

\section{Decision procedure}

For each candidate series:

\begin{enumerate}
\item \textbf{Fingerprint it} (Chapter~\ref{ch:featforensics}): ACF shape,
  lead--lag curve, cross-day continuity.
\item If it is already an SMA (triangular ramp): do not smooth ---
  \textbf{difference or deconvolve} at the detected window.
\item If its lead--lag mass is at small positive $k$ (a timing signal):
  do not smooth --- feed raw, and protect it from any pipeline stage that
  smooths globally.
\item If it is fast and noisy \emph{and} its information is slow (high
  nowcast mass spread over many lags): smooth away --- but emit the
  \emph{pair} (smooth, residual), not the smooth alone.
\item If you need denoising without delay: average across the
  cross-section, across seeds, or across models --- the delays live only
  on the time axis.
\end{enumerate}

\begin{jsbox}[: the empirical scoreboard]
Everything temporal-smoothing-related that we or the reference solutions
actually shipped, reduced to a table:
raw features (no smoothing) --- shipped;
rolling \emph{deviation} from trailing mean --- shipped, strong;
rolling std as a feature --- shipped;
market average (cross-sectional smooth) --- shipped, strong;
EWMA-30 companions --- optional extras, marginal;
temporal smoothing of predictions or labels --- never shipped by anyone,
in any solution we studied. The genre has voted.
\end{jsbox}
